\section{Checked interpolation and integrated Jensen estimates}

These are reusable analytic statements for the remaining prime-interval
argument.  The interpolation theorem has explicit holomorphy, boundedness,
and integrability hypotheses.  Its specialization to the mollifier, and
the resulting unconditional zero-density theorem, are separate steps.

\begin{definition}[Vertical square integral]
\label{density:vertical-l2}
\lean{Sidon.verticalL2,Sidon.verticalPairing}\leanok
For $\Phi:\mathbb C\to\mathbb C$ put
\[
 J(u)=\int_{\mathbb R}|\Phi(u+it)|^2\,dt.
\]
For a continuous function $g:\mathbb R\to\mathbb C$ and real $c,d$, put
\[
 P_g(z)=\int_{[c,d]}\Phi(z+it)\overline{g(t)}\,dt.
\]
The latter is an integral over a set; if $d<c$ its domain is empty.
\end{definition}

\begin{lemma}[Compact parameter integration]
\label{density:compact-parameter}
\uses{density:vertical-l2}
\lean{Sidon.continuousOn_compact_integral,Sidon.differentiableOn_compact_integral,Sidon.diffContOnCl_verticalPairing}\leanok
Suppose $a<b$ and $\Phi$ is analytic on a neighborhood of each point
of $a\leq\Re z\leq b$.  Then $P_g$ is holomorphic on
$a<\Re z<b$ and continuous on the closed strip.
\end{lemma}
\begin{proof}
For continuity, regard the closed strip as a locally compact subtype.
The integrand is jointly continuous in that parameter and in $t$.
Integration of a jointly continuous function over the compact set
$[c,d]$ is continuous: near a fixed parameter, a compact neighborhood
times $[c,d]$ supplies a uniform integrable constant bound, and
dominated convergence applies.

For differentiability at an interior point $z$, choose $r>0$ such that
the closed ball of radius $r$ about $z$ lies in the open strip.
The derivative integrand
$\Phi'(w+it)\overline{g(t)}$ is jointly continuous there.
Its norm has an upper bound $M$ on the product of that closed ball
and $[c,d]$, by compactness.  The constant $M$ is integrable on
$[c,d]$.  The integrand and its derivative are measurable, and the
pointwise complex derivative exists throughout the open ball.
The dominated differentiation theorem therefore gives
\[
 P_g'(z)=\int_{[c,d]}\Phi'(z+it)\overline{g(t)}\,dt.
\]
All neighborhoods and bounds are uniform in $t$; no interchange of an
uncontrolled infinite-domain derivative and integral is used.
\end{proof}



\begin{lemma}[Cumulative bounds imply Gaussian bounds]
\label{density:gaussian-cumulative}
\lean{Sidon.gaussianCumulativeConstant_pos,Sidon.gaussian_integral_of_cumulative}\leanok
There is an absolute constant $C>0$ with the following property.
Let $g:\mathbb R\to[0,\infty)$ be continuous, $K,N\geq0$,
$R\geq1$, and $0\leq p\leq1$.  If, for every $U\geq0$,
\[
 \int_{[-U,U]}g(t)\,dt\leq K(U+N)(1+U)^p,
\]
then $t\mapsto e^{-2t^2/R^2}g(t)$ is integrable and
\[
 \int_{\mathbb R}e^{-2t^2/R^2}g(t)\,dt
 \leq CK(R+N)R^p.
\]
\end{lemma}
\begin{proof}
Partition the real line into the measurable, pairwise disjoint sets
$A_n=\{t:\lfloor|t|/R\rfloor=n\}$ for $n\in\mathbb N$.
On $A_n$, one has $nR\leq|t|<(n+1)R$ and
$e^{-2t^2/R^2}\leq e^{-n}$: for $n\geq1$ use $2n^2\geq n$,
and for $n=0$ use nonnegativity of $t^2$.
Apply the cumulative bound with $U=(n+1)R$.  The inequalities
\[
 U+N\leq(n+2)(R+N),\qquad 1+U\leq(n+2)R,
 \qquad(n+2)^p\leq n+2
\]
give
\[
 \int_{A_n}e^{-2t^2/R^2}g(t)\,dt
 \leq K(R+N)R^p e^{-n}(n+2)^2.
\]
Each restricted integrand is integrable because $A_n$ is contained
in a compact interval and the integrand is continuous there.
The majorant is summable: expand $(n+2)^2=n^2+4n+4$ and use
summability of a polynomial times $e^{-n}$.
The summable restricted norm integrals imply global integrability.
Countable additivity over the partition and termwise comparison give
the desired estimate with
$C=1+\sum_{n\geq0}e^{-n}(n+2)^2>0$.
In particular, the constant is independent of the moving exponent $p$.
\end{proof}

\begin{theorem}[Vertical $L^2$ log-convexity]
\label{density:l2-interpolation}
\uses{density:compact-parameter}
\lean{Sidon.norm_compact_pairing_le,Sidon.norm_verticalPairing_interpolate,Sidon.verticalL2_le_interpolation,Sidon.verticalL2_logConvex}\leanok
Assume $a<b$, $a\leq\sigma\leq b$, and the preceding analyticity
hypothesis.  Assume also that $\Phi$ is bounded on the closed strip
and that $t\mapsto|\Phi(u+it)|^2$ is integrable for every
$u\in[a,b]$.  With $v=(\sigma-a)/(b-a)$,
\[
 J(\sigma)\leq J(a)^{1-v}J(b)^v.
\]
\end{theorem}
\begin{proof}
Translation invariance of Lebesgue measure gives
$\int_{\mathbb R}|\Phi(z+it)|^2dt=J(\Re z)$.
Cauchy--Schwarz and restriction of a nonnegative integral to $[c,d]$
therefore give
\[
 |P_g(z)|\leq \sqrt{J(\Re z)}
       \left(\int_{[c,d]}|g(t)|^2dt\right)^{1/2}.
\]
The boundedness of $\Phi$ and boundedness of $g$ on $[c,d]$ also
bound $P_g$ uniformly on the strip.  The scalar three-lines theorem
applies, with the holomorphy and boundary continuity proved above.

Fix $c,d$ and choose $g(t)=\Phi(\sigma+it)$, which is continuous.
Write $I=\int_{[c,d]}|g(t)|^2dt\geq0$, $A=\sqrt{J(a)}$, and
$B=\sqrt{J(b)}$.  At the real point $\sigma$, the pairing is exactly
the nonnegative real number $I$.  Three-lines yields
\[
 I\leq(A\sqrt I)^{1-v}(B\sqrt I)^v
   =A^{1-v}B^v\sqrt I.
\]
If $I=0$, the desired bound is immediate.  Otherwise divide by
$\sqrt I>0$ and square to obtain
$I\leq(A^{1-v}B^v)^2$.
This argument permits $A=0$ or $B=0$; no boundary positivity is
assumed.  Finally the increasing sets $[-L,L]$, $L\in\mathbb R$,
exhaust $\mathbb R$.  Integrability gives convergence of their
integrals to $J(\sigma)$, so the same bound holds in the limit.
Real-power multiplication and $A^2=J(a)$, $B^2=J(b)$ give the stated
formula.
\end{proof}

\begin{lemma}[Mollifier detector in Jensen's formula]
\label{density:jensen-detector}
\lean{Sidon.mollifierDetector_factor,Sidon.zeta_divisor_le_mollifierDetector,Sidon.log_norm_one_sub_sq_le,Sidon.jensen_mollifier_detector_le}\leanok
Put $f_N=\zeta M_N-1$ and $h_N=1-f_N^2$.
On a closed disc avoiding $1$, with center $c$, radius $R>0$, and
$h_N(c)\ne0$, every zeta zero is a zero of $h_N$ with at least its
zeta multiplicity, and
\[
 \sum_z\operatorname{ord}_z(h_N)
       \log\frac R{|c-z|}
 \leq \mathop{\rm av}_{|z-c|=R}|f_N(z)|^2-\log|h_N(c)|.
\]
The divisor sum is restricted to the closed disc.
\end{lemma}
\begin{proof}
The exact identity
$h_N=\zeta\,M_N(2-\zeta M_N)$ is polynomial algebra.
Both factors are analytic on the disc.  Nonvanishing of their
product at the center shows that neither factor vanishes identically
on the connected disc.  Their finite orders add under multiplication;
the companion factor has nonnegative order, proving the multiplicity
comparison.  Also
$|1-w^2|\leq1+|w|^2$ and $\log(1+x)\leq x$ for $x\geq0$ imply
$\log|1-w^2|\leq|w|^2$ when $1-w^2\ne0$.
At zero the formal real logarithm is defined as zero, so the same
inequality still holds.  Jensen's formula for the analytic detector
and monotonicity of circle averages give the assertion.
Boundary zeros are allowed by the checked closed-disc Jensen formula.
\end{proof}

\begin{lemma}[Mass of one zero under vertical integration]
\label{density:jensen-mass}
\lean{Sidon.jensen_disc_window,Sidon.jensen_log_weight_lower,Sidon.continuous_jensenVerticalWeight,Sidon.jensenVerticalWeight_integral_lower,Sidon.card_mul_jensen_mass_le_integral}\leanok
Suppose $\tfrac12\leq\sigma\leq1$, $0\leq\delta\leq\tfrac14$,
$\sigma+\delta\leq\beta\leq1$, and $4\leq\gamma\leq T$.
Put $\rho=\beta+i\gamma$, $c_y=2+iy$, $R=2-\sigma$, and
\[
 w_\rho(y)=\max\{0,\log R-\log|c_y-\rho|\}.
\]
Then
\[
 \int_3^{T+1}w_\rho(y)\,dy\geq\frac{2\delta^2}{3}.
\]
For any finite collection of such zeros, including repeated indices,
its cardinality times $2\delta^2/3$ is bounded by the integral of
the sum of these weights.
\end{lemma}
\begin{proof}
The real part of $c_y-\rho$ is $2-\beta\geq1$, so its norm is
at least one.  In particular the logarithmic weight is continuous
for every real $y$.
If $|y-\gamma|\leq\delta$, then
\[
 |c_y-\rho|^2\leq(R-\delta)^2+\delta^2
                   \leq(R-\delta/2)^2.
\]
For the second inequality the difference of right and left sides is
$R\delta-7\delta^2/4\geq0$, since $R\geq1$ and
$\delta\leq1/4$.  Both compared norms are nonnegative, so
$|c_y-\rho|\leq R-\delta/2$.
Writing this positive norm as $d$, the elementary inequality
$\log(R/d)\geq1-d/R$ gives
\[
 \log(R/d)\geq\frac\delta{2R}\geq\frac\delta3,
 \qquad R\leq3/2.
\]
The interval $[\gamma-\delta,\gamma+\delta]$ lies inside $[3,T+1]$.
Integrate the lower bound $\delta/3$ on that interval of length
$2\delta$ and use nonnegativity outside it.  Summing the resulting
inequality over a finite index set commutes with integration because
every weight is continuous on the compact interval.  Repeated indices
give exactly the multiplicity convention needed for zero counting.
\end{proof}

\begin{lemma}[Actual occurrence counts from slice bounds]
\label{density:actual-jensen-count}
\uses{density:jensen-detector,density:jensen-mass,construction:mollifier-center}
\lean{Sidon.sum_zetaZeroOccurrences,Sidon.sum_zeta_jensen_weight_le_error,Sidon.integral_circleAverage_vertical_le,Sidon.jensen_zero_occurrences_bound,Sidon.zetaZeroCount_le_low_add_high}\leanok
Let $N\geq131072$, $\tfrac12\leq\sigma_0\leq1$,
$0\leq\delta\leq\tfrac14$, and $T\geq3$.
Let $S$ be any finite subcollection of the actual positive-height zeta-zero
occurrences up to height $T$, with $\Re\rho\geq\sigma_0+\delta$
and $\Im\rho\geq4$ for each occurrence.
If, uniformly for $u\in[\sigma_0,4-\sigma_0]$,
\[
 \int_{[1,T+3]}|f_N(u+it)|^2\,dt\leq L,
\]
then
\[
 |S|\frac{2\delta^2}{3}
 \leq L+(T-2)\frac{131072}{N}.
\]
\end{lemma}
\begin{proof}
Set $R=2-\sigma_0$.  For $y\in[3,T+1]$ the disc about $2+iy$
of radius $R\leq3/2$ avoids the pole $1$: its center has imaginary
part at least three, whereas its radius is at most $3/2$.
The checked center estimate gives $h_N(2+iy)\ne0$ and
$-\log|h_N(2+iy)|\leq131072/N$.

Summing over the occurrence indices of one zero gives exactly its
analytic multiplicity times its weight.
Inside the disc, this multiplicity is the zeta divisor and is at most
the detector divisor.  The logarithmic weight is nonnegative there.
Outside the disc, the nonnegative-part weight and the restricted
detector divisor both vanish.  The detector divisor has finite support
on the compact disc, so extending a finite sum to its full divisor sum
adds only nonnegative terms.  Jensen's formula consequently bounds
the sum of occurrence weights by the error circle average plus
$131072/N$.  Restricting to $S$ can only decrease it.

Integrate in $y$.  Write the circle point as
\[
 (2+R\cos\theta)+i(y+R\sin\theta).
\]
Its real part belongs to $[\sigma_0,4-\sigma_0]$.
For fixed $\theta$, the translated height interval is contained in
$[1,T+3]$.  Nonnegativity and the slice hypothesis bound its integral
by $L$.  Joint continuity on the compact parameter rectangle supplies
absolute integrability, so Fubini exchanges the angular and vertical
integrals.  The normalized angular average of the bound $L$ is $L$.
The constant center loss integrates over an interval of length $T-2$.
The mass lower bound for every occurrence gives the conclusion.
Finally every occurrence with positive height less than four belongs
to the fixed height-four zero set.  Partitioning a finite occurrence
set at height four therefore gives
$N(\sigma,T)\leq N(0,4)+|S|$ when $S$ contains all remaining
occurrences satisfying the real-part cutoff.
\end{proof}

\begin{lemma}[Uniform optimization of the actual mollifier slices]
\label{density:optimized-slices}
\uses{construction:mollifier-gaussian,construction:mollifier-right}
\lean{Sidon.mollifier_interpolation_power_bound,Sidon.ceil_mollifier_parameters,Sidon.exists_mollifier_optimized_interior_bound,Sidon.mollifier_optimized_right_bound,Sidon.exists_mollifier_optimized_bound,Sidon.exists_optimizedMollifierSlices}\leanok
There are constants $K>0$ and $T_0$ such that for every $T\ge T_0$
and $1/2\le\sigma_0\le1$, there is an integer $N\ge131072$ with
$T\le N$ and, uniformly for $u\in[\sigma_0,4-\sigma_0]$,
\[
 \int_{[1,T+3]}|f_N(u+it)|^2\,dt
 \le K T^{(11/4)(1-\sigma_0)}(\log T)^9.
\]
One may use $N=\lceil T+3\rceil$ and
$T_0=\max\{131072,e^4\}$.
\end{lemma}
\begin{proof}
First use a variable $R\ge1$ with $\log R\ge2$, set
$N=\lceil R\rceil$ and $\delta=1/\log R$.
The ceiling inequalities give
\[
 R\le N\le2R,\qquad
 1+\log N\le1+\log(2R)\le2\log R.
\]
Here $\log2\le1$ and $\log R\ge2$ justify the last step.
The two Gaussian boundary estimates consequently give absolute
$L,B\ge1$ such that
\[
 J(1/2)\le L(\log R)R^{11/8},\qquad
 J(1+\delta)\le B(\log R)^8.
\]
Indeed $(R+N)\le3R$ on the left and $(R+N)/N\le2$ on the right.
The identity $R\,R^{3/8}=R^{11/8}$ is valid since $R>0$.

For $\sigma_0\le u\le1+\delta$, write
$\theta=(u-1/2)/(1/2+\delta)$ and $q=1-\theta$.
Then $0\le\theta,q\le1$ and
\[
 q=\frac{1+\delta-u}{1/2+\delta}
 \le2(1-\sigma_0+\delta).
\]
The numerator is nonnegative, so replacing the denominator by $1/2$
preserves the direction of the inequality.
Since $L\log R\ge1$ and $B(\log R)^8\ge1$, the powers $q,\theta$
at most one give
\begin{align*}
 J(1/2)^qJ(1+\delta)^\theta
 &\le L B(\log R)^9R^{(11/8)q}\\
 &\le LB e^{11/4}(\log R)^9R^{(11/4)(1-\sigma_0)}.
\end{align*}
For the last line use monotonicity of real powers for $R\ge1$ and
\[
 R^{(11/8)q}
 \le R^{(11/4)(1-\sigma_0)+(11/4)/\log R}
 =e^{11/4}R^{(11/4)(1-\sigma_0)}.
\]
The Gaussian recovery inequality supplies only the absolute factor
$8e^2$, so it is absorbed in the constant.

For $u\ge1+\delta$, use the direct right-half-plane mean square.
The integral on $[1,R]$ is at most that on $[0,R]$, by nonnegativity.
As $R\le N$, the latter is at most
\[
 16(\log R)^4+1024(\log R)^8+512(\log R)^5
 \le1552(\log R)^9.
\]
Since $\sigma_0\le1$ implies
$R^{(11/4)(1-\sigma_0)}\ge1$, enlarging the same constant covers
this entire half-plane.  This avoids restricting the circle's right
edge to the interpolation strip.

Finally set $R=T+3$.
For $T\ge T_0$ one has $R\le2T$ and
$\log R\le2\log T$.
With $a=(11/4)(1-\sigma_0)$, the hypotheses imply $0\le a\le11/8<2$;
therefore
\[
 R^a\le(2T)^a\le4T^a,\qquad
 (\log R)^9\le512(\log T)^9.
\]
An enlargement by $2048$ gives the claimed constant.
The integral over the closed interval equals the interval integral
because Lebesgue measure has no atoms at its endpoints.
The chosen $N$ is at least $T$ and at least $131072$.
\end{proof}

\begin{theorem}[An unconditional uniform zero-density estimate]
\label{density:uniform-zeta}
\uses{density:optimized-slices,density:actual-jensen-count,analytic:crude-zero-count}
\lean{Sidon.rpow_logarithmic_shift,Sidon.crude_rpow_le_density_scale,Sidon.jensen_center_loss_le,Sidon.exists_uniformZeroDensity_of_mollifier_slices,Sidon.exists_uniformZetaZeroDensity}\leanok
There exist $C>0$ and $U$ such that for every real $T\ge U$ and
every $0\le\sigma\le1$,
\[
 N(\sigma,T)\le C T^{(23/8)(1-\sigma)}(\log T)^{11}.
\]
Here $N(\sigma,T)$ counts the actual zeta zeros with
$0<\Im\rho\le T$ and $\Re\rho\ge\sigma$, including their actual
analytic multiplicities.  No zero-density premise is assumed.
\end{theorem}
\begin{proof}
Take the constants $K,T_0$ from the optimized slice lemma and enlarge
the threshold to $U=\max\{T_0,e^4\}$.
For $T\ge U$ put $\ell=\log T\ge4$, $\varepsilon=1/\ell\le1/4$,
and $X=T^{(23/8)(1-\sigma)}\ge1$.
The checked crude count gives a constant $D>0$ such that
$N(0,T)\le DT^{4/3}$ for $T\ge1$.

If $\sigma\le1/2+\varepsilon$, monotonicity in $\sigma$ and
\[
 \frac43\log T
 \le\frac{23}{8}+\frac{23}{8}(1-\sigma)\log T
\]
give $N(\sigma,T)\le De^{23/8}X$.
The displayed inequality follows from
$(1-\sigma)\ell\ge\ell/2-1$ and $4/3\le23/16$.
Multiplying the bound by $\ell^{11}\ge1$ proves the desired form.
This case includes every $\sigma$ below the critical line and the
small window immediately to its right.

Otherwise put $\sigma_0=\sigma-\varepsilon$.
Then $1/2\le\sigma_0\le1$ and the optimized slice lemma provides
an admissible $N$ with $T\le N$.
Let $S$ contain exactly the occurrences counted by $N(\sigma,T)$
whose heights are at least four.
Every occurrence in $S$ has real part at least
$\sigma_0+\varepsilon$.
The integrated Jensen inequality gives
\[
 |S|\frac{2\varepsilon^2}{3}
 \le K T^{(11/4)(1-\sigma_0)}\ell^9
       +(T-2)\frac{131072}{N}.
\]
The second term is at most $131072$, since $N\ge T>0$.
Also, exactly rather than asymptotically,
\[
 T^{(11/4)(1-\sigma_0)}
 =e^{11/4}T^{(11/4)(1-\sigma)}
 \le e^{11/4}X.
\]
The equality uses $\varepsilon\log T=1$; the inequality uses
$11/4\le23/8$ and $1-\sigma\ge0$.
Dividing by $2\varepsilon^2/3>0$ and enlarging $3/2$ to $2$ yields
\[
 |S|\le2\ell^2\bigl(Ke^{11/4}X\ell^9+131072\bigr).
\]
All remaining positive-height occurrences belong to the fixed finite
height-four set, of cardinality $Z=N(0,4)$.
Since $X\ell^{11}\ge1$ and $\ell^2\le X\ell^{11}$,
\[
 N(\sigma,T)
 \le\bigl(Z+2(Ke^{11/4}+131072)\bigr)X\ell^{11}.
\]
Thus the positive constant
$C=De^{23/8}+Z+2(Ke^{11/4}+131072)$ works in both cases.
\end{proof}
